\documentclass{article}
\usepackage[margin=28mm]{geometry}
\usepackage{amsmath,amssymb}
\usepackage{xurl}
\usepackage[hidelinks]{hyperref}
% Shared commands for notes in docs/paper-gaps/.

\DeclareUnicodeCharacter{2080}{\ifmmode{}_{0}\else\textsubscript{0}\fi}
\DeclareUnicodeCharacter{2081}{\ifmmode{}_{1}\else\textsubscript{1}\fi}
\DeclareUnicodeCharacter{2082}{\ifmmode{}_{2}\else\textsubscript{2}\fi}
\DeclareUnicodeCharacter{2099}{\ifmmode{}_{n}\else\textsubscript{n}\fi}

\newcommand{\Lean}{\textsc{Lean}}
\ExplSyntaxOn
\NewDocumentCommand{\leanid}{m}{
  \ifmmode
    \textsf{\detokenize{#1}}
  \else
    \group_begin:
    \urlstyle{sf}
    \tl_set:Nn \l_tmpa_tl {#1}
    \tl_replace_all:Nnn \l_tmpa_tl {\_} {_}
    \exp_args:NV \path \l_tmpa_tl
    \group_end:
  \fi
}
\ExplSyntaxOff
\providecommand{\tr}{\operatorname{tr}}
\newcommand{\tnleanIssueBase}{https://github.com/LionSR/TNLean/issues/}
\newcommand{\tnleanPrBase}{https://github.com/LionSR/TNLean/pull/}
\newcommand{\tnleanDocBase}{https://sirui-lu.com/TNLean/docs/}
\newcommand{\ghissue}[1]{\href{\tnleanIssueBase#1}{GitHub issue \##1}}
\newcommand{\ghpr}[1]{\href{\tnleanPrBase#1}{PR \##1}}
\newcommand{\doclink}[1]{\href{\tnleanDocBase#1.html}{\path{#1}}}

% Dirac notation and the identity operator, as in blueprint/src/macros/common.tex.
% \providecommand keeps note-local definitions authoritative.
\usepackage{dsfont}
\providecommand{\ket}[1]{|#1\rangle}
\providecommand{\bra}[1]{\langle #1|}
\providecommand{\braket}[2]{\langle #1 | #2 \rangle}
\providecommand{\ketbra}[2]{|#1\rangle\!\langle #2|}
\providecommand{\Id}{\mathds{1}}
\providecommand{\one}{\mathds{1}}
\providecommand{\C}{\mathbb{C}}
\providecommand{\MN}[1]{M_{#1}(\C)}
\providecommand{\MD}{M_D(\C)}

% Verdict marker: \gapnote{<kind>}{<status>}. Kind is the policy
% classification (clarification, local-correction, scope-restriction,
% unfaithful, false-source, open-gap); status is open, wip (elimination
% underway), resolved, or historical. Machine-read by the site index;
% prints a verdict line.
\newcommand{\gapnote}[2]{\par\noindent\textbf{Verdict:} #1 (#2).\par}

\providecommand{\gapnote}[2]{\par\noindent\textbf{Verdict:} #1 (#2).\par}
\title{Gaussian Moments in Hastings' Commutator Expansion}
\author{}
\date{}
\begin{document}
\maketitle
\gapnote{local-correction}{resolved}

\section{The printed coefficients}
Hastings, \emph{Solving gapped Hamiltonians locally},
\href{https://arxiv.org/abs/cond-mat/0508554v3}{arXiv:cond-mat/0508554v3},
source \path{proj3new.tex}, lines 1005--1009, expands a Gaussian-filtered
local term in iterated commutators. For even $k$ the printed coefficient
contains $(q/\Delta E^2)^k$. The normalized Gaussian in the defining
integral (source lines 194--207) has variance $q/\Delta E^2$, so the required
power is $k/2$.

\section{The expansion}
Let $q>0$, $\Delta E>0$, and $H$ be a finite-dimensional Hermitian operator.
Put $\sigma^2=q/\Delta E^2$ and $\operatorname{ad}_H(A)=HA-AH$. If $t$ has
mean-zero Gaussian distribution of variance $\sigma^2$, then
\begin{align*}
 \mathbb E(t^{2r})&=(2r-1)!!\,\sigma^{2r}, &
 \mathbb E(t^{2r+1})&=0.
\end{align*}
Expanding $e^{itH}Ae^{-itH}$ and integrating gives
\begin{align*}
 \widetilde A
 &=\sum_{r=0}^{\infty}\frac{(-1)^r}{2^r r!}
       \left(\frac q{\Delta E^2}\right)^r
       \operatorname{ad}_H^{2r}(A).
\end{align*}
Termwise integration is justified by the integrable majorant
$\|A\|e^{2\|H\||t|}$ against the Gaussian measure. Thus, in the convention
$\widetilde A=A+\sum_{k\ge1}a_k i^k\operatorname{ad}_H^k(A)$, the coefficients
are $a_k=0$ for odd $k$. For even $k\ge2$, they are
\begin{align*}
 a_k&=\frac{(k-1)!!}{k!}\left(\frac q{\Delta E^2}\right)^{k/2}.
\end{align*}
The sign $(-1)^{k/2}$ is already supplied by $i^k$. In a convention without
this factor it belongs in the coefficient instead.

\section{Scope}
The variational lower bound holds for every finite real choice of the
commutator coefficients and is unaffected by this correction. The
infinite Gaussian expansion is not asserted to be formalized by that
finite-polynomial theorem.
\end{document}
